\subsection{An ordinary peripheral obstruction at a split \texorpdfstring{$A_6$}{A6} chief step}

\begin{statement}{D.32}{Proved}{An exhausted Nielsen family obstructs the ordinary generating lift}
There is a split finite chief extension
\[
 1\longrightarrow A_6\longrightarrow\mathcal E
       \xrightarrow{\pi}\mathcal D\longrightarrow1,
 \qquad
 |\mathcal D|=4\cdot360^{12},\quad
 |\mathcal E|=4\cdot360^{13},
 \tag{D.32.1}
\]
an actual source epimorphism
$q_{\mathcal E}:\widehat F\twoheadrightarrow\mathcal E$,
and an order-two automorphism
$\alpha$ of $\mathcal D$ preserving all three
labelled peripheral conjugacy classes, such
that no epimorphism
$f:\widehat F\twoheadrightarrow\mathcal E$
satisfies
\[
 \pi f=\alpha\pi q_{\mathcal E},\qquad
                 f(x_i)\in q_{\mathcal E}(x_i)^{\mathcal E}.
 \tag{D.32.2}
\]
All base conjugators are allowed in
\textup{(D.32.2)}.
There are exactly 2160 weak lifts; each
is the graph of an explicit coordinate
projection of $\mathcal D$, conjugated
in the added $A_6$ factor, and has image
a complement to that factor.
For the original base map, the corresponding
ordinary peripheral fibre has 3240 weak
lifts, of which 1080 are epimorphic and
2160 have complementary image.

The maximal pro-five quotients of
$\mathcal E$ and $\mathcal D$ are trivial.
For every prescribed finite-five-group
multiplier-one peripheral automorphism
$\beta$, taking the actual products with
its source quotient gives both an ordinary
obstructed base shadow $\alpha\times\beta$
and an explicitly admitted base shadow
$1\times\beta$, with the same induced
pro-five shadow $\beta$.
\end{statement}

\begin{proof}
Use the actual group
$G=\operatorname{Aut}(A_6)$ of D.31,
its inner subgroup $S=\operatorname{Inn}(A_6)$,
and its quotient $O=G/S=\mathbf F_2^2$.
The elements of $G$ below are specified
by their images of the same generators
$a=(0123)(45)$ and $b=(01234)$.
Write $[i_0i_1i_2i_3i_4i_5]$ for the
permutation $j\mapsto i_j$.
Define $U\in G$ by
\[
                  U(a)=[103452],\qquad
                  U(b)=[023451].
 \tag{D.32.3}
\]
It has order six.  The following table
defines fifteen actual automorphisms
$V_{H_j}$ and $V_{L_j}$:
\[
\begin{array}{c|cc@{\qquad}c|cc}
 &a\text{ image}&b\text{ image}
 & &a\text{ image}&b\text{ image}\\\hline
 H_0&[103524]&[325140]&L_0&[123054]&[510423]\\
 H_1&[104532]&[423015]&L_1&[123054]&[542103]\\
 H_2&[130254]&[025413]&L_2&[140523]&[513402]\\
 H_3&[130254]&[352410]&L_3&[150432]&[314250]\\
 H_4&[153204]&[214530]&L_4&[203154]&[524301]\\
 H_5&[230514]&[153042]&L_5&[320154]&[054213]\\
 H_6&[240531]&[134205]&&&\\
 H_7&[241503]&[054132]&&&\\
 H_8&[401523]&[035214]&&&
\end{array}
 \tag{D.32.4}
\]
For each label $T$, also use $T$ to
denote the actual homomorphism
\[
 T:\widehat F\twoheadrightarrow G,\qquad
       (T(x_1),T(x_2),T(x_3))
          =(U,V_T,(UV_T)^{-1}).
 \tag{D.32.5}
\]
All fifteen maps are onto, and all
have the same outer map
\[
 \theta:\widehat F\twoheadrightarrow O,\qquad
       (x_1,x_2,x_3)\longmapsto(e_1,e_2,e_1+e_2).
 \tag{D.32.6}
\]
These assertions are verified by
literal automorphism multiplication
and enumeration of the generated
1440-element subgroup, using the
complete automorphism construction
in D.31.

\paragraph{The two complete Nielsen families.}
Let $\mathcal U$ be the $S$-class
of $U$, let $\mathcal V$ be the
$S$-class of $V_{H_0}$, and put
\[
 \mathcal W_H=((UV_{H_0})^{-1})^S,\qquad
 \mathcal W_L=U\mathcal W_HU^{-1}.
 \tag{D.32.7}
\]
Their sizes are $120,90,90,90$.
The classes $\mathcal W_H$ and
$\mathcal W_L$ are distinct.
The exact tuple counts and orbit
decompositions are
\[
\begin{array}{c|ccc}
 \text{profile}&\text{literal triples}
     &S\text{-conjugacy orbits}&\text{generated order}\\\hline
 (\mathcal U,\mathcal V,\mathcal W_H)&3240&9&1440\\
 (\mathcal U,\mathcal V,\mathcal W_L)&2160&6&1440.
\end{array}
 \tag{D.32.8}
\]
The representatives are precisely
the nine $H_j$ and the six $L_j$
in \textup{(D.32.4)--(D.32.5)}.
For completeness, the enumeration
fixes the first element $U$ and
uses its inner centralizer of order
three on the remaining admissible
second elements.  There are 27
admissible second elements for
$\mathcal W_H$ and eighteen for
$\mathcal W_L$; their orbits all
have size three.  Every representative
generates $G$.  Conjugating the
first element through its 120
values gives the literal counts
in \textup{(D.32.8)}.  Thus these
counts retain properness and the
complete orbit decomposition.

Put $K=\ker\theta$.  Each of the
fifteen maps sends $K$ onto $S$.
Their fifteen kernels on $K$ are
pairwise distinct.  This last
statement is verified by actual
relators, not inferred only from
the orbit labels: for every pair
$T,T'$ the certificate saves a
word $w_{T,T'}\in K$ with
$T(w_{T,T'})=1$ and
$T'(w_{T,T'})\ne1$.
There are 105 such words.
For example
\[
 w=x_2x_1x_2^2x_1^2x_2^{-1}x_1^{-1}
 \tag{D.32.9}
\]
has $H_0(w)=1$, whereas $H_1(w)$
sends $a$ to $[320154]$ and
$b$ to $[324105]$.
Its two exponent parities are
zero, so it lies in $K$.
All 105 saved words are checked
both in their two actual images
and in \textup{(D.32.6)}.

This also gives the precise
comparison between kernel orbits
and the relevant automorphisms.
If two of the $K$-kernels agreed,
their onto maps $K\to S$ would
differ by some
$\eta\in\operatorname{Aut}(S)=G$.
For any $g\in\widehat F$, the
conjugation actions of $T'(g)$
and $\eta T(g)\eta^{-1}$ on
that onto image would then agree.
Their difference centralizes $S$
in $G$ and is trivial.
Hence $T'=\eta T\eta^{-1}$ as
actual maps on all of $\widehat F$.
The literal relators rule this out,
including automorphisms outside
the inner subgroup $S$.

\paragraph{The actual common-outer products and their source.}
Define
\[
 \begin{split}
 \mathcal D
   &=\{(g_1,\ldots,g_{12})\in G^{12}:
                  \overline g_1=\cdots=\overline g_{12}\in O\},\\
 \mathcal E
   &=\{(g_0,g_1,\ldots,g_{12})\in G^{13}:
                  \overline g_0=\cdots=\overline g_{12}\in O\},
 \end{split}
 \tag{D.32.10}
\]
with coordinatewise multiplication.
Projection forgetting $g_0$ is
$\pi$ in \textup{(D.32.1)}.
Its kernel is the actual first
copy of $S$; the map
$d\mapsto(d_1,d)$ is a splitting,
where $d_1$ means the first
coordinate of $d\in\mathcal D$.
Thus the extension is split and
its kernel is simple.

The source maps, in the exact
coordinate order used below, are
\[
 \begin{split}
 q_{\mathcal D}
   &=(H_0,L_0,H_1,L_1,\ldots,H_5,L_5),\\
 q_{\mathcal E}
   &=(H_6,H_0,L_0,H_1,L_1,\ldots,H_5,L_5).
 \end{split}
 \tag{D.32.11}
\]
They land in \textup{(D.32.10)}
because of their common outer
map.  They are epimorphisms.
Indeed their coordinate maps on
$K$ are onto the nonabelian
simple group $S$ and have
pairwise distinct kernels by
the relators above.
The simple subdirect-product
argument therefore gives all
of $S^{12}$ or $S^{13}$.
Here is the exact implication:
induct on the number of
coordinates; if the last
coordinate were constrained
after projection to the preceding
full product, it would be an
onto homomorphism from that
product to $S$.  The commuting
normal images of its simple
factors force this homomorphism
to factor through one coordinate,
repeating a kernel.  This
contradicts the saved witness.
Together with the common onto
map to $O$, full generation
and the two orders in
\textup{(D.32.1)} follow.

Let $\alpha$ swap each adjacent
pair of coordinates of
$\mathcal D$.
This is an actual automorphism
of the displayed group, of
order two, and fixes its
common quotient $O$.
It preserves each labelled
peripheral conjugacy class.
For the first peripheral all
coordinates equal $U$ and
no conjugator is needed.
For the second, choose
$c_j\in S$ with
\[
 c_j V_{H_j}c_j^{-1}=V_{L_j}.
 \tag{D.32.12}
\]
Then the tuple
$(c_0,c_0^{-1},\ldots,c_5,c_5^{-1})$
belongs to $\mathcal D$ and
performs all six swaps.
For the third peripheral,
choose $b_j\in G$ of outer
class $e_1$ sending the
$H_j$ third value to the
$L_j$ third value.
These elements exist by
\textup{(D.32.7)}, and the
tuple $(b_0,b_0^{-1},\ldots,b_5,b_5^{-1})$
again belongs to $\mathcal D$:
all twelve outer coordinates
are $e_1$.
The certificate supplies every
$c_j,b_j$ as an actual saved
automorphism and verifies
each conjugacy equation.
Denote the three resulting
base conjugators by $h_i$.
Their common outer coordinates
are respectively
\[
                           0,\quad0,\quad e_1.
 \tag{D.32.13}
\]

\paragraph{The complete ordinary centralizer union.}
Write $d_i=q_{\mathcal D}(x_i)$
and let $\rho:\mathcal D\to O$
be the common-outer projection.
The exact centralizer images are
\[
 \rho(C_{\mathcal D}(d_1))=\langle e_1\rangle,\qquad
 \rho(C_{\mathcal D}(d_2))=\langle e_2\rangle,\qquad
 \rho(C_{\mathcal D}(d_3))=\langle e_1+e_2\rangle.
 \tag{D.32.14}
\]
To prove this, the image of
the centralizer of any relevant
coordinate in $G$ is exactly
the indicated line.  Equivalently,
that line is the stabilizer in
$O$ of its $S$-conjugacy class.
This is computed for the actual
classes in \textup{(D.32.7)}.
The image of a centralizer
in the common-outer product
is the intersection of these
coordinate images.  Equality
with the intersection follows
by choosing in each coordinate
a centralizing element with
the same prescribed outer value;
the resulting tuple belongs
to \textup{(D.32.10)}.
The certificate saves such
a tuple for both values of
every line in
\textup{(D.32.14)}.

Every base conjugator sending
$d_i$ to $\alpha(d_i)$ belongs
to the actual coset
$h_iC_{\mathcal D}(d_i)$.
Thus all their possible outer
values, allowing the full
ordinary conjugacy condition,
are
\[
 \langle e_1\rangle,\qquad
 \langle e_2\rangle,\qquad
 e_1+\langle e_1+e_2\rangle
                       =\{e_1,e_2\}.
 \tag{D.32.15}
\]
The first line stabilizes
$\mathcal U$, the second
stabilizes $\mathcal V$, and
both elements of the third
coset send $\mathcal W_H$
to $\mathcal W_L$.
Consequently the entire
ordinary union of allowable
first-coordinate triples over
$\alpha q_{\mathcal D}$ is
exactly
\[
                   (\mathcal U,\mathcal V,\mathcal W_L).
 \tag{D.32.16}
\]
This is also sufficient,
not just necessary.
For any triple in that profile,
choose first-coordinate
conjugators of outer values
$0,0,e_1$.  Pair them with
the actual twelve-coordinate
$h_i$ above; the pairs belong
to $\mathcal E$ and realize
the three required conjugacies.
The product relation is
coordinatewise.  Hence the
2160 triples in the second
row of \textup{(D.32.8)}
give exactly all ordinary
weak lifts in
\textup{(D.32.2)}.

\paragraph{Exhaustion and the exact two fibres.}
Each triple in
\textup{(D.32.16)} is uniquely
an $S$-conjugate of one of
the six $L_j$ representatives.
All six are already present
in the base source.
More exactly, under the
swapped base map
$\alpha q_{\mathcal D}$,
the coordinate in the
position formerly occupied
by $H_j$ equals $L_j$.
If the added first-coordinate
map is $cL_jc^{-1}$, its
image in $\mathcal E$ is
therefore the graph of the
actual homomorphism
\[
 \mathcal D\longrightarrow\mathcal E,\qquad
       d\longmapsto
             (c\,\operatorname{pr}_{H_j}(d)c^{-1},d).
 \tag{D.32.17}
\]
Here $c\in S$ and
$\operatorname{pr}_{H_j}$
means that fixed coordinate
position, not the old source
map.
This graph is a complement
to the first $S$ factor and
has order $|\mathcal D|$.
It is a proper subgroup
of $\mathcal E$ of index
360.  The six indices and
360 distinct inner conjugates
give the 2160 weak lifts,
all complementary and none
epimorphic.  This proves
the ordinary obstruction.

For the original base map,
the same calculation using
\textup{(D.32.14)} gives the
entire first-coordinate profile
$(\mathcal U,\mathcal V,\mathcal W_H)$.
Of its nine complete kernel
classes, six $H_0,\ldots,H_5$
are present in the base and
give 2160 graph lifts.
The three remaining classes
$H_6,H_7,H_8$ have distinct
$K$-kernels from every base
coordinate by the actual
relators.  Each therefore
gives a full subdirect
product and 360 epimorphisms.
This proves the original
counts $3240,2160,1080$ and
also verifies the original
source choice $H_6$ directly.

\paragraph{The exact pro-five comparison.}
Both groups in
\textup{(D.32.10)} have a
perfect normal subgroup,
$S^{12}$ or $S^{13}$, and
quotient $O$ of order four.
Every map to a finite
five-group is therefore
trivial, by the same
explicit quotient argument
as in D.31.
For any finite five-group
source quotient $L$ with
a multiplier-one peripheral
automorphism $\beta$, the
paired sources generate
$\mathcal E\times L$ and
$\mathcal D\times L$:
the subdirect-product
comparison would otherwise
give a nontrivial common
five-group quotient.
Their maximal pro-five
maps are the actual
coordinate projections
to $L$.
The automorphism
$\alpha\times\beta$ is
peripheral on the base,
using the pairs of the
conjugators above and the
prescribed conjugators
for $\beta$.
Any ordinary epimorphic
lift would project to a
lift in
\textup{(D.32.2)}, which
has just been excluded.
For $1\times\beta$ the
explicit paired map
$(q_{\mathcal E},\beta q_L)$
is an epimorphic peripheral
lift.  Both base
automorphisms induce the
same $\beta$ on the
maximal pro-five quotient.

The computation files in
the \texttt{code} directory
are
\texttt{a\_peripheral\_a6\_nielsen.py}
and
\texttt{a\_peripheral\_a6\_ordinary.py},
with executed outputs and
complete JSON certificates.
The first retains every
fixed-first orbit and its
generated subgroup.
The second retains all
105 literal kernel witnesses,
all twelve-coordinate
conjugators, the complete
centralizer-image vectors,
and the six exact graph
coordinate maps.
The common-outer product
orders are derived from
the proved simple-factor
generation argument; they
are not reported as an
enumeration of those large
groups.
The final status is
\texttt{PASS\_A6\_ORDINARY\_CHIEF\_OBSTRUCTION}.

This is an ordinary finite
chief-step obstruction,
with all base-centralizer
choices included.
The explicit products and
pro-five projections show
its precise limitation for
D.1--D.2: it excludes a
finite branch while the
same prescribed pro-five
shadow has the admitted
branch displayed above.
\end{proof}
